\documentclass{article}
% ---------------------------------------------------------------
% Póster de revisión de literatura (A0 vertical)
% Finanzas I (055D), UNCP, 2026-II. Compilar con pdfLaTeX.
% ---------------------------------------------------------------
\usepackage[paperwidth=841mm,paperheight=1189mm,margin=0mm]{geometry}
\usepackage[utf8]{inputenc}
\usepackage[T1]{fontenc}
\IfFileExists{spanish.ldf}{\usepackage[spanish,es-noshorthands,es-nodecimaldot]{babel}}{}
\usepackage{helvet}
\renewcommand{\familydefault}{\sfdefault}
\usepackage[letterspace=120]{microtype}
\usepackage{anyfontsize}
\usepackage{graphicx}
\usepackage{xcolor}
\usepackage[most]{tcolorbox}
\usepackage{enumitem}
\usepackage{multicol}
\usepackage{xurl}
\usepackage[hidelinks]{hyperref}
\usepackage{qrcode}
\urlstyle{same}
\pagestyle{empty}
\setlength{\parindent}{0pt}
\hyphenpenalty=10000 \exhyphenpenalty=10000

% ---------------- Paleta ----------------
\definecolor{navy}{HTML}{8E0C11}
\definecolor{navyLight}{HTML}{B8121A}
\definecolor{gold}{HTML}{00A3E0}
\definecolor{acento}{HTML}{A6E1F7}
\definecolor{goldSoft}{HTML}{E8F6FC}
\definecolor{teal}{HTML}{353435}
\definecolor{tealSoft}{HTML}{F1EEEE}
\definecolor{ink}{HTML}{231F20}
\definecolor{muted}{HTML}{5E5757}
\definecolor{line}{HTML}{E5D7D7}
\definecolor{page}{HTML}{F6F1F1}
\definecolor{mist}{HTML}{F7D3D5}
\pagecolor{page}

% ---------------- Tamaños (ajustados por póster) ----------------
\newcommand{\bodysize}{\fontsize{32}{42.2}\selectfont}
\newcommand{\cardtitlesize}{\fontsize{37.8}{45.4}\selectfont}
\newcommand{\smallsize}{\fontsize{26.2}{34.1}\selectfont}

% ---------------- Estilos de tarjeta ----------------
\tcbset{
  card/.style={enhanced, sharp corners, colback=white, colframe=line, boxrule=0.9mm,
    left=14mm, right=14mm, top=12mm, bottom=12mm, fontupper=\bodysize\color{ink}\raggedright,
    overlay={\fill[#1] (frame.north west) rectangle ([yshift=-5mm]frame.north east);}},
  card/.default=navy,
}
\newcommand{\cardtitle}[2][navy]{{\color{#1}\cardtitlesize\bfseries\textls[110]{\MakeUppercase{#2}}}\par\vspace{7mm}}
\newcommand{\bul}{\raisebox{0.18em}{\color{gold}\rule{0.5em}{0.5em}}}
\setlist[itemize]{label=\bul, leftmargin=1.25em, labelsep=0.6em, itemsep=0.55em, topsep=0pt, parsep=0pt}
\newcommand{\pill}[3]{\tcbox[on line, enhanced, arc=7mm, boxrule=0.8mm, colback=#1, colframe=#2,
  left=7mm, right=7mm, top=3mm, bottom=3mm, fontupper=\fontsize{26}{30}\selectfont\bfseries]{#3}}

\begin{document}

% ======================= ENCABEZADO =======================
\begin{tcolorbox}[enhanced, sharp corners, colback=navy, boxrule=0pt, width=\paperwidth,
  left=34mm, right=34mm, top=20mm, bottom=20mm,
  overlay={\fill[gold] (frame.south west) rectangle ([yshift=5mm]frame.south east);
           \fill[navyLight] ([xshift=-160mm]frame.north east) -- (frame.north east) -- ([yshift=-160mm]frame.north east) -- cycle;}]
\begin{minipage}[c]{92mm}\centering\includegraphics[width=\linewidth]{Escudo_UNCP.png}\end{minipage}\hfill
\begin{minipage}[c]{555mm}\centering
{\color{acento}\fontsize{25}{30}\selectfont\bfseries\textls[160]{REVISIÓN DE LITERATURA \quad|\quad FINANZAS I \quad|\quad TEMA 61}\par}
\vspace{10mm}
{\color{white}\fontsize{66}{74.6}\selectfont\bfseries The effects of short selling on price discovery: A study for Borsa Istanbul\par}
\vspace{8mm}
{\color{mist}\fontsize{38}{46.4}\selectfont\itshape Efectos de las ventas en corto sobre el descubrimiento de precios: un estudio para Borsa Istanbul\par}
\vspace{10mm}
{\color{white}\fontsize{29}{35}\selectfont\bfseries Baki Cem Şahin y Fatih Kuz\par}
\vspace{9mm}
\pill{white}{white}{\color{navy}Q1}\hspace{6mm}\pill{navy}{white}{\color{white}\mdseries \textit{Borsa Istanbul Review}, 21(2), 133-138 (2021)}
\end{minipage}\hfill
\begin{minipage}[c]{92mm}\centering\includegraphics[width=\linewidth]{Logo_FE.png}\end{minipage}
\end{tcolorbox}
\vspace{-1mm}
\begin{tcolorbox}[enhanced, sharp corners, colback=goldSoft, boxrule=0pt, width=\paperwidth,
  left=34mm, right=34mm, top=8mm, bottom=8mm, halign=center]
{\color{navy}\fontsize{27}{36}\selectfont \textbf{Póster elaborado por:} Meglinho Sneyder Rojas Poma \enspace\textbar\enspace Código 2024200522B\par \fontsize{24}{32}\selectfont Escuela Profesional de Economía, Universidad Nacional del Centro del Perú \enspace\textbar\enspace Finanzas I (055D), 2026-II\par}
\end{tcolorbox}

\vfill
% ======================= CUERPO =======================
\hspace*{34mm}\begin{minipage}{773mm}
\begin{tcolorbox}[enhanced, sharp corners, colback=navy, boxrule=0pt, left=16mm, right=16mm, top=12mm, bottom=11mm,
  fontupper=\bodysize\raggedright, overlay={\fill[gold] (frame.north west) rectangle ([yshift=-5mm]frame.north east);}]
\begin{minipage}[t]{363.5mm}\raggedright{\color{acento}\smallsize\bfseries\textls[110]{\MakeUppercase{Mercado}}}\par\vspace{2mm}{\color{white}Borsa Istanbul (Turquía)}\end{minipage}\hfill\begin{minipage}[t]{363.5mm}\raggedright{\color{acento}\smallsize\bfseries\textls[110]{\MakeUppercase{Método}}}\par\vspace{2mm}{\color{white}Regresión de eficiencia informativa}\end{minipage}
\end{tcolorbox}
\end{minipage}

\vfill
\hspace*{34mm}\begin{minipage}{773mm}
\begin{tcbraster}[raster columns=3, raster equal height, raster column skip=16mm]
\begin{tcolorbox}[card]\cardtitle{Contexto}Las ventas en corto son polémicas: en las crisis algunos reguladores las prohíben porque las consideran manipuladoras, mientras otros estudios dicen que aceleran el ajuste de los precios.\end{tcolorbox}
\begin{tcolorbox}[card]\cardtitle{Problema}No se sabía si las ventas en corto mejoran o empeoran la formación de precios en Borsa Istanbul, sobre todo en los días de caídas extremas.\end{tcolorbox}
\begin{tcolorbox}[card]\cardtitle{Objetivo}Analizar el efecto de las ventas en corto sobre la eficiencia informativa de Borsa Istanbul y el comportamiento de los vendedores en corto en días de caídas extremas.\end{tcolorbox}
\end{tcbraster}
\end{minipage}

\vfill
\hspace*{34mm}\begin{minipage}{773mm}
\begin{tcolorbox}[card]\cardtitle{Metodología}
\begin{minipage}[t]{0.48\linewidth}
\begin{itemize}
\item Se construye una medida de eficiencia informativa de los precios.
\item Se regresa sobre la proporción de ventas en corto y variables de control.
\item Se analiza la conducta de los vendedores alrededor de días con caídas extremas.
\end{itemize}
\end{minipage}\hfill
\begin{minipage}[t]{0.48\linewidth}
\begin{itemize}
\item Las reversiones se clasifican en pequeñas, grandes y de sobrerreacción según cuánto se recupera el precio.
\item Se controla por volatilidad anormal, volumen anormal, tamaño y otras variables.
\end{itemize}
\end{minipage}
\end{tcolorbox}
\end{minipage}

\vfill
\hspace*{34mm}\begin{minipage}{773mm}
\begin{tcolorbox}[card]\cardtitle{Figura clave del estudio}
\begin{minipage}[c]{291.1mm}\centering
\fcolorbox{line}{white}{\includegraphics[width=\dimexpr291.1mm-2\fboxsep-2\fboxrule\relax]{figura.png}}
\end{minipage}\hfill
\begin{minipage}[c]{\dimexpr\linewidth-291.1mm-18mm\relax}\raggedright
{\color{navy}\bfseries Figura 1.}\enspace{\itshape Ventas en corto alrededor de fechas con retornos negativos extremos: retorno (línea continua) y ratio de ventas en corto (línea punteada, eje derecho), del día $-20$ al $+20$. Paneles: continuación, reversión pequeña, reversión grande y sobrerreacción.}\par\vspace{8mm}
{\color{teal}\smallsize\bfseries\textls[110]{CÓMO LEERLA}}\par\vspace{4mm}
\begin{itemize}
\item En todos los paneles hay una caída fuerte del retorno el día 0, seguida de un retorno positivo en las reversiones.
\item El ratio de ventas en corto sube en los días previos a la caída extrema.
\item El día de la caída las ventas en corto bajan: los vendedores en corto no agravan la caída.
\end{itemize}
\vspace{8mm}
{\color{muted}\smallsize Figura tomada de Şahin y Kuz (2021).}
\end{minipage}
\end{tcolorbox}
\end{minipage}

\vfill
\hspace*{34mm}\begin{minipage}{773mm}
\begin{tcolorbox}[card=teal]\cardtitle[teal]{Resultados}
\begin{tcolorbox}[enhanced, sharp corners, colback=tealSoft, boxrule=0pt, left=14mm, right=14mm, top=10mm, bottom=11mm,
  borderline west={4mm}{0mm}{teal}, fontupper=\bodysize\color{ink}\raggedright]
{\color{teal}\smallsize\bfseries\textls[110]{HALLAZGO CENTRAL}}\par\vspace{4mm}
{\bfseries Las ventas en corto mejoran la eficiencia informativa de los precios.}
\end{tcolorbox}
\vspace{11mm}
\begin{minipage}[t]{0.48\linewidth}
\begin{itemize}
\item Las ventas en corto aumentan antes de las caídas extremas.
\item No agravan las caídas: disminuyen el día de la caída extrema.
\end{itemize}
\end{minipage}\hfill
\begin{minipage}[t]{0.48\linewidth}
\begin{itemize}
\item Su volumen crece cuanto mayor es la recuperación posterior.
\item El coeficiente de las ventas en corto anormales crece con el tamaño de la reversión.
\end{itemize}
\end{minipage}
\end{tcolorbox}
\end{minipage}

\vfill
\hspace*{34mm}\begin{minipage}{773mm}
\begin{tcbraster}[raster columns=3, raster equal height, raster column skip=16mm]
\begin{tcolorbox}[card, raster multicolumn=2]\cardtitle{Conclusiones}Prohibir las ventas en corto en las crisis puede quitarle al mercado un mecanismo que ayuda a que los precios reflejen la información.\end{tcolorbox}
\begin{tcolorbox}[card=gold, colback=goldSoft]\cardtitle{Aporte al tema 61}Muestra que los mecanismos que ayudan a formar precios importan. Cuando los precios se ajustan con lentitud, como en un mercado poco negociado, las betas estimadas pueden salir sesgadas.\end{tcolorbox}
\end{tcbraster}
\end{minipage}

\vfill
% ======================= PIE =======================
\begin{tcolorbox}[enhanced, sharp corners, colback=navy, boxrule=0pt, width=\paperwidth,
  left=34mm, right=34mm, top=14mm, bottom=11mm,
  overlay={\fill[gold] (frame.north west) rectangle ([yshift=-5mm]frame.north east);}]
\begin{minipage}[c]{245mm}\raggedright
{\color{acento}\smallsize\bfseries\textls[110]{FUENTE DE DATOS DEL ESTUDIO}}\par\vspace{5mm}
{\color{white}\smallsize Datos de ventas en corto y precios de las acciones de Borsa Istanbul.}
\end{minipage}\hfill
\begin{minipage}[c]{385mm}\raggedright
{\color{acento}\smallsize\bfseries\textls[110]{REFERENCIA}}\par\vspace{5mm}
{\color{white}\smallsize\hangindent=1.2em\hangafter=1 Şahin, B. C., y Kuz, F. (2021). The effects of short selling on price discovery: A study for Borsa Istanbul. \textit{Borsa Istanbul Review}, \textit{21}(2), 133-138. \url{https://doi.org/10.1016/j.bir.2020.09.004}\par}
\end{minipage}\hfill
\begin{minipage}[c]{100mm}\centering
\tcbox[enhanced, sharp corners, colback=white, boxrule=0pt, left=3mm, right=3mm, top=3mm, bottom=3mm]{\qrcode[hyperlink, height=46mm]{https://doi.org/10.1016/j.bir.2020.09.004}}\par\vspace{4mm}
{\color{acento}\fontsize{17}{21}\selectfont\bfseries\textls[100]{ESCANEA Y LEE EL ARTÍCULO}}
\end{minipage}
\end{tcolorbox}
\end{document}
